\chapter{The cold start and the erasure bridge}
\label{chap:coldstart}

The three preceding chapters describe an \emph{idealised} erasure: a relation \Erases{} between
Lean terms and \lbox{} terms, a pass relation \Lower{} that abstracts fixpoint blocks and prunes
eliminators, and a \lbox{} operational semantics. None of them mentions the program that actually
runs when a user types \code{\#erase foo to "foo.ast"}. That program is \code{Erasure.erase}, a
\code{CoreM} computation over a mutable \code{Erasure.ErasureState} --- a constant registry, an
inductive registry, the emitted \lbox{} global declarations \code{gdecls} and an inlining list ---
driven by eighteen mutually recursive \code{EraseM} functions and punctuated by \code{panic!}
fall-throughs, \code{Lean.Meta} queries and compiler preprocessing passes. This chapter is about
\emph{that} program. It says what a run \textbf{is} (a decomposition down to two run equations from
the empty state), what every run \textbf{produces} (an output-shape induction with no hypotheses at
all), what a run \textbf{maintains} (a registry invariant relating a fixed specification environment
to the state), and what the verification still \textbf{assumes} about the connection between the two.

The name ``cold start'' refers to the empty initial state. The hinge is
Theorem~\ref{thm:erase_run_ok}: a successful \code{Erasure.erase e cfg} unfolds, by pure
monad-transformer algebra and with no hypotheses, into a preparation run followed by a
\code{visitExpr} run, both at the \emph{default} \code{ErasureState} and at the entry reader. Every
invariant in the chapter is engineered to be free at that state. The centre of gravity is a single
induction, Theorem~\ref{thm:visitExpr_shapeW}, run once through the erasure family's fixpoint
induction principle over all eighteen members: instead of proving one fact at a time about the
eraser, the development abstracts ``what a predicate of the erasure state must be closed under''
into a sixteen-field interface and spends the induction on whichever predicate a consumer needs.
Its unconditional instance, Theorem~\ref{thm:visitExpr_shape_all}, is the strongest
hypothesis-free statement anywhere in the development.

The chapter's second half is the relation-facing side. \code{RegInvShape'}
(Definition~\ref{def:RegInvShape-prime}) is MetaCoq's environment-erasure judgement recast as a loop
invariant of a stateful registration pass, and a preservation kit tiles the registration path step by
step; four projections turn it into the environment-side results the capstone needs, and
Theorem~\ref{thm:bridgeEnv_of_regInv} assembles them into exactly the two fields of
\code{ErasureBridge}. The chapter closes by stating plainly where the chain stops: nothing in the
repository produces a \code{RegInvShape'} for a state the shipping eraser actually reaches. The
reader should see two finished layers --- the entry point and the induction, both spent by the
capstone and by the refinement of Chapter~\ref{chap:refinement} --- plus one prepared but
unconnected layer. The files are \code{LeanToLambdaBox/ColdStartRun.lean},
\code{ColdStartInduction.lean}, \code{ColdStartShape.lean} and \code{Bridge.lean}, together with the
four \code{RegInvShape'} projections of \code{SpecEnv.lean} and the three bridge declarations of
\code{Capstone.lean}.

\section{The entry point: what a run of the shipping eraser is}
\label{sec:coldstart:entry}

\begin{theorem}[The cold-start entry point]
\label{thm:erase_run_ok}
\lean{LeanToLambdaBox.erase_run_ok}
\leanok
\uses{def:Erasure-erase, def:Erasure-EraseM, def:Erasure-prepare_erasure, def:Erasure-visitExpr}
A successful run of the shipping entry point \emph{is} a preparation run followed by a term walk,
both from the empty state. Precisely: if \code{Erasure.erase e cfg} succeeds at the \code{CoreM}
world $w$, returning the program $p$ and the inlining list \code{inls}, then there are a prepared
term \code{pe}, a \lbox{} term $t$, states \code{sp}, \code{sf} and worlds such that
\code{prepare\_erasure e \{\} \{ config := cfg \}} succeeds at $(\code{pe}, \code{sp})$,
\code{visitExpr pe sp \{ config := cfg \}} succeeds at $(t, \code{sf})$, and
$p = \code{.untyped sf.gdecls (some t)}$ with $\code{inls} = \code{sf.inlinings}$. Nothing is
assumed: no condition on $e$, none on the configuration, none on the state.
\end{theorem}
\begin{proof}
\leanok
\uses{lem:Erasure-run_bind_ok, lem:Erasure-run_pure, def:Erasure-EraseM}
Unfold \code{Erasure.erase}, which binds \code{Erasure.run} and returns the program built from the
final state. A \code{CoreM}-level bind inversion splits the outer bind; the file's own
\code{run\_eq} rewrites \code{Erasure.run x cfg} to \code{x \{\} \{ config := cfg \}} by
\code{rfl}, which is where the \emph{default} state and the entry reader --- empty local context, no
fixvar map, no declaration universe parameters --- become visible; the run algebra's bind inversion
then splits the inner \code{EraseM} bind into the preparation and the walk, and the tail is a pure
return. The \code{CoreM} helpers \code{core\_bind\_ok} and \code{core\_pure} used here are stated in
the same file and are not given separate nodes.
\end{proof}
\begin{remark}
This is the first step of the capstone's proof and the reason every invariant of this chapter is
designed to cost nothing at entry.
\end{remark}
\begin{theorem}[The preparation pass, with the \code{csimp} gate off]
\label{thm:run_prepare_erasure_ok}
\lean{LeanToLambdaBox.run_prepare_erasure_ok}
\leanok
\uses{def:Erasure-prepare_erasure, def:ConfigPinned, def:Erasure-EraseM}
If the reader's configuration has $\code{csimp} = \code{false}$, then a successful
\code{prepare\_erasure} run leaves the erasure state untouched and is exactly its four
\code{CoreM}-lifted transforms, in order: \code{Erasure.replaceUnsafeRecNames},
\code{Lean.Compiler.LCNF.macroInline}, \code{Lean.Compiler.LCNF.inlineMatchers}, and
\code{macroInline} again. The four sub-runs are handed back individually, which is where a per-pass
evaluation-preservation premise can be read. The gate $\code{csimp} = \code{false}$ is the first
conjunct of \code{ConfigPinned}.
\end{theorem}
\begin{proof}
\leanok
\uses{lem:Erasure-run_bind_ok, lem:Erasure-run_liftCoreM_ok, lem:Erasure-run_pure}
Unfold \code{prepare\_erasure} and step every monadic bind. All four transforms are \code{CoreM}
actions reached through \code{liftM}, so the run algebra's state-transparency lemma reports that
none of them touches \code{ErasureState}. The single branch that is not a plain lift --- a
\code{Lean.Core.transform} at \code{EraseM} through \code{MonadControlT}, guarded by
\code{@[csimp]} --- is dead under the hypothesis, which is why the gate is a premise and not a
conclusion.
\end{proof}
\begin{lemma}[The preparation leaves the state alone]
\label{lem:run_prepare_erasure_state}
\lean{LeanToLambdaBox.run_prepare_erasure_state}
\leanok
\uses{def:Erasure-prepare_erasure, def:ConfigPinned}
At $\code{csimp} = \code{false}$, a successful \code{prepare\_erasure} run ends in the state it
started in.
\end{lemma}
\begin{proof}
\leanok
\uses{thm:run_prepare_erasure_ok}
The first projection of Theorem~\ref{thm:run_prepare_erasure_ok}.
\end{proof}
\begin{remark}
This is consumed inside \code{erasure\_bridge\_of\_run} (Theorem~\ref{thm:erasure_bridge_of_run}),
and again --- as a \emph{parameter} rather than a citation --- by
Lemma~\ref{lem:runClosedW_runConcl}, because \code{ColdStartRun.lean} sits \emph{above}
\code{ColdStartInduction.lean} in the import graph. State transparency here is exactly what lets
Theorem~\ref{thm:erase_run_ok}'s \code{visitExpr} run be read at the \textbf{empty} state rather
than at an unknown post-pass state.
\end{remark}
\begin{theorem}[The preparation preserves the source evaluation]
\label{thm:prepare_sound}
\lean{LeanToLambdaBox.prepare_sound}
\leanok
\uses{def:Erasure-prepare_erasure, def:EraserAsks, def:SEval, def:SEvalFlags, def:ConfigPinned, def:mkApps}
Assume the eraser's own obligation bundle \code{EraserAsks} and $\code{csimp} = \code{false}$. If
\code{prepare\_erasure} takes $e$ to \code{pe}, then for every application spine \code{args}, every
compiler body table, every level scope, every flag point, every model local context and every value
$v$: if the source evaluation \SEval{} takes $\mkApps\ e\ \code{args}$ to $v$, it takes
$\mkApps\ \code{pe}\ \code{args}$ to $v$ as well.
\end{theorem}
\begin{proof}
\leanok
\uses{thm:run_prepare_erasure_ok, asm:EraserAsks-passes_sound}
Theorem~\ref{thm:run_prepare_erasure_ok} exposes the four sub-runs; each of the three distinct
passes is shown by simplification to be a member of \code{preparePasses}, and the field
\code{EraserAsks.passes\_sound} is applied four times, chaining the four evaluations. The spine is
quantified rather than pushed through, because each pass is a whole-tree \code{Lean.Core.transform}
walk, so $f(\mkApps\ e\ \code{args})$ is \emph{not} $\mkApps\ (f\ e)\ \code{args}$; the capstone
reads its observable at a spine, so the spine must be carried through the passes rather than
reconstructed after them.
\end{proof}
\begin{remark}
This is one of the two places the capstone touches this chapter. It is complete \emph{modulo an
assumption}: \code{EraserAsks.passes\_sound} is a hypothesis field owned by this repository, not a
theorem, so the semantic adequacy of the preprocessing is assumed, not proved. The module docstring
of \code{ColdStartRun.lean} claims of the whole file that there is ``no erasure relation, no
lean4lean, no assumption''; that sentence is false of this theorem.
\end{remark}

\section{What every run produces: one induction, eighteen motives}
\label{sec:coldstart:induction}

\begin{definition}[Closure of a state predicate under the eraser's steps]
\label{def:RunClosed}
\lean{LeanToLambdaBox.RunClosed}
\leanok
\uses{def:Erasure-EraseM, def:Erasure-register_inductive, def:Erasure-prepare_erasure, def:Erasure-nonrecConstState, def:NoFix, def:LBClosed}
A \code{Prop}-valued structure with six fields, one per place the eighteen-member family touches the
erasure state: \code{inl} (the \code{@[inline]} bookkeeping cons), \code{ax} and \code{reg} (the two
registration primitives, stated at a \emph{run}, so that \code{addAxiom}'s \code{panic!}
fall-through and \code{register\_inductive}'s state-changing cold branch are both covered),
\code{prep} (the preprocessing pass), and \code{nrc}/\code{rc} (the two registering exits of
\code{Erasure.visitMutual}). \code{nrc} is allowed to \emph{assume} the shape of the term being
stored --- fix-free, closed at level $0$, applied form; \code{rc} is allowed to assume only
closedness and applied form of each fixpoint projection, since the projection is itself a
\code{fix} node and so is never fix-free to begin with. These facts are false of an arbitrary
term and are proved by the same induction that consumes them.
\end{definition}
\begin{remark}
\code{RunClosed} is world-blind and therefore unable to carry a configuration condition; it is
inhabited only at the trivial predicate and survives exactly to give the unconditional output-shape
corollary, Theorem~\ref{thm:visitExpr_shape_all}.
\end{remark}
\begin{definition}[The per-call conclusions]
\label{def:ShapeCW}
\lean{LeanToLambdaBox.ShapeCW, LeanToLambdaBox.ShapeC}
\leanok
\uses{def:NoFix, def:LBClosed, def:Erasure-EraseM}
$\code{ShapeC}\ Q\ s\ s'\ t$ says: from $Q$ at the entry state, $Q$ at the exit state and the
produced \lbox{} term is fix-free, de Bruijn closed at level $0$, and in applied (spine) constructor
form. $\code{ShapeCW}\ \code{Cfg}\ P\ \code{ctx}\ s\ w\ s'\ w'\ t$ is its world-indexed twin: from
$P$ at the entry state-and-world and \code{Cfg} of the reader's configuration, $P$ at the exit
state-and-world and the same three conjuncts about the term.
\end{definition}
\begin{remark}
\code{NoBlock} is not a boxing condition; it forbids exactly one node shape, a constructor node
carrying a non-empty block argument list. The eraser has one constructor construction site,
\code{Erasure.visitConstructor}, nullary by design, so the induction \emph{concludes} applied form
rather than assuming it --- the point at which this development departs from MetaCoq, which instead
ships a verified block/applied transformation. The configuration premise deliberately sits
\emph{after} the run equation in the motive, which is the shape the family's admissibility lemmas
recognise for fixpoint induction.
\end{remark}
\begin{definition}[Closure of a world-indexed state predicate]
\label{def:RunClosedW}
\lean{LeanToLambdaBox.RunClosedW}
\leanok
\uses{def:RunClosed, def:Erasure-isErasable, def:Erasure-register_inductive, def:Erasure-prepare_erasure, def:Erasure-EraseM}
A \code{Prop}-valued structure with \textbf{sixteen} fields, parameterised by a condition \code{Cfg}
on the reader's configuration and by a predicate $P$ of the state \emph{and} the opaque
\code{IO.RealWorld} token. Ten fields --- \code{oracle}, \code{inferType}, \code{constInfo},
\code{getEnv}, \code{logInfo}, \code{isInstance}, \code{fresh}, \code{declInfo}, \code{ctorArity},
\code{casesInfo} --- are the ambient primitives the family calls (the relevance oracle
\code{Erasure.isErasable}, \code{Lean.Meta.inferType}, \code{Lean.getConstInfo},
\code{Lean.getEnv}, \code{Lean.logInfo}, \code{Lean.Meta.isInstance}, \code{Lean.mkFreshFVarId},
\code{Lean.Compiler.LCNF.getDeclInfo?}, \code{Lean.Compiler.LCNF.getCtorArity?},
\code{Lean.getCasesInfo?}); each needs its own clause because it advances the world even though it
leaves the state alone. The remaining six are \code{RunClosed}'s, with two differences: \code{prep}
and \code{reg} additionally read \code{Cfg} of the reader's configuration, and \code{reg} carries
the \emph{provenance} of the inductive value it is called at --- either a \code{getConstInfo} run
returned it, or the configuration's \code{nat} field is \code{.machine}, the two machine-numeral
arms of \code{Erasure.visitCases}.
\end{definition}
\begin{remark}
The world index is forced by two facts the bridge needs that no state predicate can express: a
\emph{relation} between two states (\code{Erasure.RunConcl}) and a \emph{bound on the name
generator}, which reads the world. The two configuration-reading clauses are forced by
\code{prepare\_erasure}'s \code{@[csimp]} branch --- a \code{Lean.Core.transform} at \code{EraseM}
that no \code{liftM} lemma reaches --- and by \code{register\_inductive}'s pruning mask. The ten
primitive clauses are one for one the calls specified by the \code{ErasureSpec} fields
\code{prim\_monotone}, \code{lookup\_adequate} and \code{fresh\_names}.
\end{remark}
\begin{lemma}[A state predicate is a world-blind world-indexed one]
\label{lem:runClosedW_of_runClosed}
\lean{LeanToLambdaBox.runClosedW_of_runClosed}
\leanok
\uses{def:RunClosed, def:RunClosedW}
If $Q$ is closed under the eraser's steps as a state predicate, then $\lambda s\,\_.\ Q\ s$ is
closed under them as a world-indexed one, at any configuration condition.
\end{lemma}
\begin{proof}
\leanok
\uses{def:RunClosed, def:RunClosedW, lem:Erasure-run_liftCoreM_ok}
The ten primitive fields are discharged by the state-transparency lemmas of the run algebra --- each
primitive's run leaves \code{ErasureState} unchanged, so $Q$ survives by rewriting. The six
remaining fields are \code{RunClosed}'s own, with \code{Cfg} simply unread.
\end{proof}
\begin{remark}
This is why there is exactly \emph{one} induction over the erasure family in the file: the
state-only induction is a corollary of the world-indexed one, not a second proof.
\end{remark}
\begin{lemma}[The stored fixpoint block is closed and in applied form]
\label{lem:rec_block_closed}
\lean{LeanToLambdaBox.rec_block_closed, LeanToLambdaBox.rec_block_noBlock}
\leanok
\uses{def:LBClosed, def:NoFix, def:toBvar, def:Erasure-mkLambda}
Suppose a block's definition list \code{defs} has one entry per member, the block's
fixpoint-variable name list has the same length, and every entry's body is the closure --- the fold
of \code{toBvar} over the reversed, indexed name list --- of some de Bruijn closed \lbox{} term.
Then the $j$-th projection of the fixpoint node is closed at level $0$ for every $j$; and, with the
closedness hypothesis replaced by applied form, that projection is in applied form for every $j$.
\end{lemma}
\begin{proof}
\leanok
\uses{lem:lbClosed_fix_of_bodies, def:LBClosed}
Closedness of a fixpoint node unfolds to ``every body is closed at the block's length'', and each
abstraction step over the block's names raises the level by one, so as many abstractions as the
block has members carry a level-$0$ body to a body closed at that level. The applied-form sibling
has no binder arithmetic: applied form of a fixpoint node is applied form of each body, discharged
directly through the fold. The premise on the fixvar name list is the shipping code's
\code{fixvarnames := names.map remove\_unsafe\_rec}, \ie{} a length-of-map rewrite.
\end{proof}
\begin{lemma}[The binder helpers, world-indexed]
\label{lem:run_withLocalDecl_okW}
\lean{LeanToLambdaBox.run_withLocalDecl_okW, LeanToLambdaBox.run_withLocalDef_okW, LeanToLambdaBox.run_lambdaMonocular_okW, LeanToLambdaBox.run_letMonocular_okW, LeanToLambdaBox.run_forallMonocular_okW, LeanToLambdaBox.run_lambdaMonocularOrIntro_okW, LeanToLambdaBox.run_lambdaOrIntroToArity_okW}
\leanok
\uses{def:Erasure-mkLambda, def:Erasure-EraseM}
Seven decompositions of the eraser's binder-opening helpers, each reporting two facts about the
continuation the helper calls: its reader has the \textbf{same configuration} as the caller's (only
the local context, the fixvar map and the level scope move), and the predicate $P$ survives to the
continuation's entry --- the freshness run that advanced the world included. The monocular family
additionally carries the \code{panic!} fall-through as an explicit disjunct, so panics are
discharged rather than excluded. The last helper peels exactly \code{arity} binders by structural
recursion on the arity.
\end{lemma}
\begin{proof}
\leanok
\uses{lem:Erasure-run_panic, lem:Erasure-run_pure, lem:Erasure-run_bind_ok}
Each helper is unfolded and its binds stepped. The run algebra's own versions hand the continuation
back at an \emph{unconstrained} reader and drop the fresh-identifier run; a world-indexed induction
needs both, because the two configuration-reading clauses of \code{RunClosedW} are evaluated at the
continuation's reader and because the world moved.
\end{proof}
\begin{remark}
The arity-peeling helper's panic disjunct reports $P$ at the exit state and world rather than
equality of worlds, since that fall-through returns at the world the binders opened so far --- an
asymmetry with its four siblings, documented in the source.
\end{remark}
\begin{lemma}[The two registering exits, world-indexed]
\label{lem:run_nonrec_exit_okW}
\lean{LeanToLambdaBox.run_nonrec_exit_okW, LeanToLambdaBox.run_rec_exit_okW}
\leanok
\uses{def:RunClosedW, def:Erasure-visitMutual, def:Erasure-nonrecConstState}
Two decompositions of \code{Erasure.visitMutual}'s registering exits, stated so that the
world-indexed predicate survives them. For the non-recursive exit: given the \code{RunClosedW}
clauses it exercises, an abstract induction hypothesis for the body erasure, and a reader
transformation that preserves the configuration, $P$ holds at the exit. For the block exit: the
same, at the \emph{two} reader updates the block loop makes (one for the fresh fixvar identifiers,
one per sibling for its own \code{levelParams}), walking the block's members with a loop invariant
that accumulates, per sibling, closedness and applied form of the open body together with its
closure equation.
\end{lemma}
\begin{proof}
\leanok
\uses{lem:rec_block_closed, lem:Erasure-run_bind_ok, lem:Erasure-run_list_mapM_ok, thm:Erasure-run_mkDef_ok, lem:Erasure-run_pure}
Both are the run algebra's exit lemmas with the sub-run's reader pinned to the caller's
configuration through the hypothesis that the reader transformation preserves the configuration,
discharged by \code{rfl} at the call site. The block exit's loop invariant is exactly the input
Lemma~\ref{lem:rec_block_closed} needs and the input the \code{rc} clause of \code{RunClosedW}
consumes.
\end{proof}
\begin{lemma}[The two numeral dispatches, enumerated]
\label{lem:visitCases_match_tri}
\lean{LeanToLambdaBox.visitCases_match_tri, LeanToLambdaBox.visitConstructor_match_quad}
\leanok
\uses{def:Erasure-visitExpr, def:ConfigPinned}
The dispatch on the pair (type name, \code{nat} mode) inside \code{Erasure.visitCases} is a
trichotomy --- the machine-\code{Nat} arm, the machine-\code{Int} arm (both of which force the
\code{nat} mode to be \code{.machine}) and the generic arm --- and the dispatch on the pair
(\code{nat} mode, constructor name) of \code{Erasure.visitConstructor} is a four-way case:
machine-\code{Nat.zero}, machine-\code{Nat.succ}, and two generic arms.
\end{lemma}
\begin{proof}
\leanok
Split on the elaborator-generated matcher, then close each arm by \code{rfl}. The lemmas are stated
against the generated matchers rather than against the shipping definitions, because a
two-discriminant match with \code{Name} \emph{literal} patterns cannot be taken apart at a
hypothesis whose subject is the match \emph{applied} to the monad's five arguments.
\end{proof}
\begin{remark}
These two gate the machine-numeral arms of the induction, which is where a
\code{register\_inductive} call for \code{Bool} happens with a provenance that is not a
\code{getConstInfo} run --- the second disjunct of \code{RunClosedW}'s \code{reg} clause. The
matcher \emph{indices} are load-bearing: editing either shipping match renumbers them and breaks the
proof, loudly, as a build error. The matcher names are deliberately absent from the declarations
cited above.
\end{remark}
\begin{theorem}[The world-indexed induction, all eighteen motives]
\label{thm:visitExpr_shapeW}
\lean{LeanToLambdaBox.visitExpr_shapeW}
\leanok
\uses{def:RunClosedW, def:ShapeCW, def:Erasure-visitExpr, def:Erasure-visitMutual}
Let \code{Cfg} be a condition on configurations and $P$ a predicate of the erasure state and the
world, closed under the eraser's steps in the sense of Definition~\ref{def:RunClosedW}. Then at
every reader whose configuration satisfies \code{Cfg}, every successful run of any of the eighteen
members of the erasure family --- \code{visitExpr}, \code{visitLiteral}, \code{visitConstructor},
\code{visitConst}, \code{get\_constant\_kername}, \code{visitMutual}, \code{visitAppArgs},
\code{visitLet}, \code{visitLambda}, \code{visitProj}, \code{visitApp}, \code{visitConstApp},
\code{visitCtorEta}, \code{visitCtorEtaGo}, \code{visitCasesEta}, \code{visitCasesEtaGo},
\code{visitCases}, \code{visitAlt} --- carries $P$ from its entry state-and-world to its exit
state-and-world; and every member that returns a \lbox{} term returns one that is fix-free, de
Bruijn closed and in applied form.
\end{theorem}
\begin{proof}
\leanok
\uses{lem:visitCases_match_tri, lem:run_withLocalDecl_okW, lem:run_nonrec_exit_okW, lem:rec_block_closed, lem:Erasure-run_panic, lem:Erasure-run_bind_ok, def:NoFix, def:LBClosed}
One application of the erasure family's mutual fixpoint induction principle, with the eighteen
admissibility obligations discharged by the family's admissibility lemmas, followed by eighteen step
cases. Each step unfolds the shipping member, steps every monadic bind, feeds each ambient primitive
to the matching field of \code{RunClosedW}, opens binders through the world-indexed binder lemmas,
and discharges every \code{panic!} arm through the run algebra's panic lemma --- a panic
\emph{succeeds} and returns the default term, whose shape is fix-free, closed and applied. The two
registering exits go through Lemma~\ref{lem:run_nonrec_exit_okW} and the two numeral dispatches
through Lemma~\ref{lem:visitCases_match_tri}.
\end{proof}
\begin{remark}
Three motives deviate from the uniform \code{ShapeCW} shape: the \code{visitAppArgs} motive
additionally \emph{takes} the accumulator seed's three conjuncts as hypotheses; the \code{visitAlt}
motive concludes closedness of the alternative's body under its own binders; and the
\code{get\_constant\_kername} and \code{visitMutual} motives return no term, so $P$'s survival is
all they say. The proof needs \code{set\_option maxHeartbeats 4000000}, worth recording for build
reproducibility. Its measured footprint contains no non-standard axiom at all: pure run
algebra, no lean4lean, no \code{sorryAx}.
\end{remark}
\begin{theorem}[The state-only induction, and its non-vacuity]
\label{thm:visitExpr_shape}
\lean{LeanToLambdaBox.visitExpr_shape, LeanToLambdaBox.runClosed_true}
\leanok
\uses{def:RunClosed, def:ShapeCW, def:Erasure-visitExpr}
From $\code{RunClosed}\ Q$: every successful run of a member of the erasure family preserves $Q$,
and every member returning a \lbox{} term returns a fix-free, closed, applied-form one. Moreover the
trivial predicate is \code{RunClosed}, so the premise is satisfiable and the theorem is not
vacuously true.
\end{theorem}
\begin{proof}
\leanok
\uses{thm:visitExpr_shapeW, lem:runClosedW_of_runClosed, def:RunClosed}
Instantiate Theorem~\ref{thm:visitExpr_shapeW} at the trivial configuration condition and at
$\lambda s\,\_.\ Q\ s$ through Lemma~\ref{lem:runClosedW_of_runClosed}, feeding each of the eighteen
conjuncts a trivial proof for the configuration premise. Non-vacuity is immediate: at the trivial
predicate every field's conclusion is \code{True}.
\end{proof}
\begin{theorem}[The output shape of the shipping eraser, unconditionally]
\label{thm:visitExpr_shape_all}
\lean{LeanToLambdaBox.visitExpr_shape_all, LeanToLambdaBox.visitExpr_noFix_closed, LeanToLambdaBox.visitExpr_noBlock, LeanToLambdaBox.visitExpr_output_shape}
\leanok
\uses{def:NoFix, def:LBClosed, def:Erasure-visitExpr}
Every successful run of \code{Erasure.visitExpr} returns a \lbox{} term $t$ that is fix-free,
closed at level $0$ and in applied form: no fixpoint node in spine position, no loose de Bruijn
index, and no constructor node carrying block arguments. There are \textbf{no hypotheses at all}
--- none on the state, none on the source expression, none on the configuration --- and the eraser's
sixteen panic sites (four \code{panic!} calls and twelve \code{unreachable!} calls, per
\code{doc/panics.md}) are discharged rather than excluded. Three thin corollaries project
the first two conjuncts, the third conjunct alone, and the generic packaging.
\end{theorem}
\begin{proof}
\leanok
\uses{thm:visitExpr_shape}
The first motive of Theorem~\ref{thm:visitExpr_shape} at the trivial state predicate, whose state
half is discarded and whose output half is left standing with no hypotheses.
\end{proof}
\begin{remark}
This is the chapter's headline unconditional theorem and the strongest hypothesis-free statement in
the development. It is the term half of \code{Output.LBWfPeregrine}'s closedness conjunct, the
non-fix premise of the fix-unfolding metatheory, and --- through the fix-free-and-closed corollary
--- the closedness input to Lemma~\ref{lem:RegInvShape-prime-constCons}. The repository's
\code{doc/panics.md} names it as the metatheory that covers the panic sites, and it proves outright
the former hypothesis \code{VisitExprRunConcl}, since deleted. The two remaining corollaries have no
consumer in the repository.
\end{remark}
\begin{lemma}[Conjunction of closed predicates]
\label{lem:RunClosedW-and}
\lean{LeanToLambdaBox.RunClosedW.and}
\leanok
\uses{def:RunClosedW}
Two world-indexed predicates closed under the eraser's steps at the same configuration condition are
closed jointly.
\end{lemma}
\begin{proof}
\leanok
\uses{def:RunClosedW}
Field by field: each of the sixteen clauses is the pair of the corresponding clauses of the two
hypotheses.
\end{proof}
\begin{remark}
This is the composition operator. It is what lets the state half, the generator half and --- in
Chapter~\ref{chap:refinement} --- the registry half of one run conclusion be established by a
\emph{single} pass of the induction instead of three.
\end{remark}
\begin{lemma}[The state half of a run conclusion]
\label{lem:runClosedW_runConcl}
\lean{LeanToLambdaBox.runClosedW_runConcl}
\leanok
\uses{def:RunClosedW, def:Erasure-RunConcl, def:Erasure-StateLe}
Fix a base state $s_0$. Given that the preparation pass leaves the state alone at every reader
satisfying \code{Cfg}, the predicate ``the state has only grown canonically since $s_0$'' ---
\code{Erasure.RunConcl}, which packages the \code{Erasure.StateLe} order on the three registries and
on \code{gdecls} together with a canonicity clause on the constant registry --- is closed under
every step of the erasure family.
\end{lemma}
\begin{proof}
\leanok
\uses{thm:Erasure-run_addAxiom_ok, thm:Erasure-run_register_inductive_runConcl, lem:Erasure-run_liftCoreM_ok, lem:Erasure-runConcl_nonrecConstState, def:Erasure-RunConcl}
The ten primitive clauses are state-transparency rewrites. Each of the four writers is a
\code{RunConcl} step already proved in the run algebra, composed with the incoming \code{RunConcl}
by transitivity; \code{addAxiom}'s post-state, its \code{panic!} fall-through included, comes from
the run algebra's \code{addAxiom} lemma.
\end{proof}
\begin{remark}
The preparation hypothesis is a \emph{parameter} rather than a citation because
Lemma~\ref{lem:run_prepare_erasure_state} lives in a module \emph{above} this one in the import
graph; it is discharged at the call sites.
\end{remark}
\begin{lemma}[The generator half of a run conclusion]
\label{lem:runClosedW_gen}
\lean{LeanToLambdaBox.runClosedW_gen}
\leanok
\uses{def:RunClosedW, def:ConfigPinned, def:ErasureSpec}
Given the primitive specifications \code{ErasureSpec}, a base world $w_0$, and bounds saying that
\code{prepare\_erasure} and \code{register\_inductive} only advance the name generator, the
predicate ``the \code{Lean.NameGenerator} has only advanced since $w_0$'' is closed under every step
of the erasure family at a pinned configuration.
\end{lemma}
\begin{proof}
\leanok
\uses{asm:ErasureSpec-prim_monotone, asm:ErasureSpec-lookup_adequate, asm:ErasureSpec-fresh_names, asm:ErasureSpec-oracle_refl}
Each ambient primitive's bound is the \code{ErasureSpec} clause that specifies it; the four state
writers leave the world alone. The configuration condition is pinned to \code{ConfigPinned} here
rather than left abstract, because the two parameters are stated at that condition.
\end{proof}
\begin{remark}
The two bounds taken as parameters are Lemma~\ref{lem:run_register_inductive_gen} and the
preparation's own generator bound, both proved, both declared below this module in the import graph.
The same file also contains \code{visitExpr\_runConcl\_gen}, which composes this lemma with
Lemma~\ref{lem:runClosedW_runConcl} through Lemma~\ref{lem:RunClosedW-and} and spends the pair at
Theorem~\ref{thm:visitExpr_shapeW}; it has \emph{no consumer in the repository}, the live version
being \code{visitExpr\_runConcl} of Chapter~\ref{chap:refinement}, which performs the same
composition and additionally carries the registry-model conjunct. It is given no node here.
\end{remark}
\begin{lemma}[Closedness of the emitted bodies]
\label{lem:runClosedW_closedBodies}
\lean{LeanToLambdaBox.runClosedW_closedBodies, LeanToLambdaBox.visitExpr_closedBodies, LeanToLambdaBox.closedBodies_empty}
\leanok
\uses{def:RunClosedW, def:ElimDecl}
Given that the preparation pass leaves the state alone, ``every body the run has emitted is de
Bruijn closed'' is closed under every step of the erasure family; hence the emitted environment of a
\code{visitExpr} run has closed bodies whenever it did at entry, and at the cold entry it does,
vacuously.
\end{lemma}
\begin{proof}
\leanok
\uses{thm:visitExpr_shape_all, thm:Erasure-run_register_inductive_cold_ok, def:Erasure-ConstExt}
The ambient primitives and the preparation leave \code{gdecls} alone; \code{addAxiom} and
\code{register\_inductive} add only body-less entries and inductive blocks, neither of which
declares a body, the cold branch's prefix being read off the run algebra's constant-extension
report; and \code{visitMutual}'s two exits store exactly the terms whose closedness the induction
itself proves.
\end{proof}
\begin{remark}
Closedness of the emitted bodies is the one clause of \code{RegInvShape'} mentioning no
specification environment, hence a fact about the run alone. The doc comment calls it the
\emph{environment} half of \code{Output.LBWfPeregrine}'s closedness conjunct, but that role is in
fact filled at every rung by a different mechanism --- the capstone's own well-formedness binder,
decided by the Boolean checker of Theorem~\ref{thm:lbWfPeregrine_of_check}. All three declarations
of this node have no in-repository consumer.
\end{remark}

\section{Constructor saturation, on the run}
\label{sec:coldstart:ctorsat}

\begin{lemma}[One argument out per argument in]
\label{lem:run_visitAppArgs_mkApps}
\lean{LeanToLambdaBox.run_visitAppArgs_mkApps}
\leanok
\uses{def:Erasure-visitExpr, def:LBTerm-mkApps}
Every successful \code{visitAppArgs f as} run returns its seed applied to exactly one \lbox{}
argument per source argument: the result is $\mkApps\ f\ \code{ts}$ for some list \code{ts} whose
length is the size of \code{as}.
\end{lemma}
\begin{proof}
\leanok
\uses{lem:Erasure-run_list_forIn_ok-prime, lem:LBTerm-mkApps_concat}
Induction on the array fold \code{visitAppArgs} performs, with the spine-extension lemma turning
each step's application into an extension of the argument list.
\end{proof}
\begin{remark}
\code{visitAppArgs} is the eraser's single application-spine builder, hence the single place a
spine's length is fixed. This is the length half of the saturation argument.
\end{remark}
\begin{definition}[The saturation statement and its mask hypothesis]
\label{def:CtorSatConcl}
\lean{LeanToLambdaBox.CtorSatConcl, LeanToLambdaBox.MaskKeeps}
\leanok
\uses{def:Erasure-register_inductive, def:Erasure-EraseM, def:LBTerm-mkApps, def:LBTerm}
$\code{CtorSatConcl}\ \code{ctx}\ \code{cctx}\ \code{ref}\ c\ s\ n\ t$ says that the emitted term
$t$ is either the panic value or the \lbox{} constructor node of the constructor the environment
reports for $c$, in \textbf{applied (spine) form with an empty block argument list}, applied to
exactly $n$ arguments whenever the source occurrence carried at least the constructor's parameters
plus fields. $\code{MaskKeeps}\ s\ \code{ctx}\ c$ is the hypothesis shape it is proved under: for
every run of the three queries the eraser makes at $c$ --- the constant's info, its inductive
type's info, then \code{register\_inductive} --- the argument mask reported for $c$ keeps every
field.
\end{definition}
\begin{remark}
\code{MaskKeeps} is stated over the run's own three queries rather than over the configuration
alone, because with constructor-argument pruning off \code{register\_inductive} computes an
all-keep mask, but a state that \emph{already} holds the block reports its cached masks instead.
Nothing in the repository proves an instance of \code{MaskKeeps}, and the definition's own doc
comment concedes the premise is stronger than a consumer would need; it is moreover evaluated at the
\emph{entry} state of the enclosing \code{visitExpr} run.
\end{remark}
\begin{theorem}[A saturated constructor occurrence is emitted saturated]
\label{thm:visitExpr_ctorSat}
\lean{LeanToLambdaBox.visitExpr_ctorSat, LeanToLambdaBox.visitApp_ctorSat, LeanToLambdaBox.visitConstructor_ctorSat}
\leanok
\uses{def:CtorSatConcl, def:ConfigPinned, def:Erasure-visitExpr}
Suppose the configuration pins $\code{nat} = \code{.peano}$ and
$\code{extern} = \code{.preferLogical}$, the mask hypothesis holds at $c$, the source term's head is
the constant $c$, the environment answers for $c$ as a constructor of arity \code{arity} and not as
a \code{casesOn}, and the occurrence carries at least \code{arity} arguments. Then a successful
\code{visitExpr} run on it satisfies \code{CtorSatConcl} at the occurrence's own argument count: the
emitted term is the constructor node applied to exactly as many arguments as the occurrence carried,
or the panic value.
\end{theorem}
\begin{proof}
\leanok
\uses{def:CtorSatConcl, lem:run_visitAppArgs_mkApps, def:Erasure-register_inductive, def:Erasure-isErasable, lem:Erasure-run_panic}
Three layers. At the construction site, \code{Erasure.visitConstructor} emits the constructor node
applied to three slices --- the parameter slice, the mask-selected field slice and the
over-application tail; under a total mask the middle slice becomes a truncation, and two
arithmetic lemmas recompose the three slices into the whole argument array, so the emitted spine
carries exactly as many arguments as the occurrence. The \code{@[extern]} exit is excluded by the
\code{extern} pin and both machine-\code{Nat} exits by the \code{nat} pin; every other exit is a
\code{panic!} and lands in the default disjunct. At the dispatch, a constant-headed application
whose head is a constructor of that arity and not a \code{casesOn}, visited at \code{arity}
arguments or more, reaches \code{visitConstructor} at exactly its own argument array: the
\code{casesOn} exit is excluded by the eliminator hypothesis, the plain-constant exit by the arity
hypothesis, the eta-expansion exit by saturation. At the top, \code{visitExpr} dispatches to
\code{visitApp}, with the erasability gate absorbed by the default disjunct --- an erasable
occurrence returns $\Box$, not a spine.
\end{proof}
\begin{remark}
This is the applied-versus-block constructor convention established \emph{on the run} rather than by
a transformation on the output; MetaCoq instead ships a verified block/applied transformation
([S], Sec. 7.4). A non-vacuity fixture in the same file shows that at a two-argument constructor the
emitted spine has exactly two arguments and is not the panic value, so the saturated disjunct is
genuinely reachable. The whole layer --- \code{MaskKeeps}, \code{CtorSatConcl}, the three theorems
and the fixture --- has no consumer anywhere in the repository; \code{doc/rework/01-DESIGN.md}
describes this theorem as ``lifted through \code{RegInvShape'}'', and that lift does not exist and
cannot exist until the open item of Section~\ref{sec:coldstart:bridge} is closed.
\end{remark}

\section{Decomposing a mutual-block visit}
\label{sec:coldstart:decomp}

\begin{definition}[The inlining slack]
\label{def:InlineExt}
\lean{LeanToLambdaBox.InlineExt}
\leanok
\uses{def:Erasure-EraseM, def:Erasure-StateLe, def:Erasure-RunConcl}
$\code{InlineExt}\ s\ s'$ says $s'$ differs from $s$ at most in the \code{inlinings} field: the
constant registry, the inductive registry and the emitted declarations all agree. It is reflexive,
transitive, holds across one \code{inlinings} cons, and implies both \code{Erasure.StateLe} and
\code{Erasure.RunConcl}.
\end{definition}
\begin{remark}
\code{Erasure.visitMutual}'s \code{@[inline]} bookkeeping prefix and tail write to the
\code{inlinings} field and to nothing else, and \code{Erasure.erase} returns that field
\emph{separately} from the program, so it is erasure-irrelevant --- but a decomposition of a
\code{visitMutual} call still has to say so. This is why the ``fourth exit'' of \code{visitMutual}
appears as slack inside each disjunct of Lemma~\ref{lem:run_visitMutual_decomp} rather than as a
case of its own.
\end{remark}
\begin{lemma}[The state effect of one mutual-block visit]
\label{lem:run_visitMutual_decomp}
\lean{LeanToLambdaBox.run_visitMutual_decomp, LeanToLambdaBox.run_inline_tail_decomp, LeanToLambdaBox.run_inline_prefix_decomp, LeanToLambdaBox.run_nonrec_exit_decomp, LeanToLambdaBox.run_rec_exit_decomp, LeanToLambdaBox.run_rec_exit_siblings, LeanToLambdaBox.run_rec_exit_siblings_closed}
\leanok
\uses{def:InlineExt, def:Erasure-visitMutual, def:Erasure-addAxiomState, def:Erasure-nonrecConstState}
A successful \code{visitMutual n} run first performs a declaration-info query at the \code{CoreM}
layer, and then takes one of three registering exits, modulo \code{InlineExt} slack at the entry:
the axiom exit; the non-recursive constant exit --- which hands back the inner
\code{prepare\_erasure} and \code{visitExpr} runs on the declaration's own value, under the reader
with no fixvar map and the declaration's own level parameters, followed by the non-recursive
constant state; or the recursive block exit. Supporting lemmas expose the per-sibling runs of the
block exit, together with each sibling's naming and closure equations, and record that each
sibling's \emph{open} body is fix-free and de Bruijn closed.
\end{lemma}
\begin{proof}
\leanok
\uses{def:InlineExt, thm:visitExpr_shape_all, thm:Erasure-run_addAxiom_ok, lem:Erasure-run_liftCoreM_ok, lem:Erasure-run_list_mapM_ok, thm:Erasure-run_mkDef_ok}
\code{visitMutual} is unfolded and its binds stepped; the inlining prefix and tail are peeled by
their own decomposition lemmas into \code{InlineExt} slack; each of the three exits is then
decomposed separately, the block exit's member loop being walked with an existentially loaded loop
invariant that keeps the per-sibling \emph{runs} rather than a state predicate's worth of their
consequences. Closedness of the per-sibling open bodies is Theorem~\ref{thm:visitExpr_shape_all},
which has no hypotheses.
\end{proof}
\begin{remark}
The lemma deliberately does \emph{not} decide which exit was taken: that depends on opaque runtime
data, so a consumer either pins the declaration or handles all three. What it buys over a state
predicate is the \emph{identification} of the dependency: the declaration-info run itself, the
preparation's input pinned to the declaration's own value, and the callee's reader pinned so that
only the fixvar map and the level parameters move. The proof needs
\code{set\_option maxHeartbeats 1000000}. This whole layer --- roughly half of
\code{ColdStartRun.lean} --- has no consumer anywhere in the repository: the live delta path goes
through Chapter~\ref{chap:refinement}'s environment steps instead, and
\code{doc/rework/06-REPAIRS-W4.md} records the intended consumer as not landed. It is machinery
built for an argument that has not been made.
\end{remark}

\section{The cold-start registry invariant}
\label{sec:coldstart:reginv}

\begin{definition}[An inductive block, emitted unchanged]
\label{def:IndEmitted}
\lean{LeanToLambdaBox.IndEmitted}
\leanok
\uses{def:IndInfo, def:LBTerm-envLookup, def:GlobalDeclarations}
For a source type former $n$, $\code{IndEmitted}\ \code{env}\ \Gamma_{\mathrm{spec}}\ \Gamma\ n$
says that whatever mutual-inductive block the specification environment $\Gamma_{\mathrm{spec}}$
declares at $n$'s block kername is declared, \textbf{unchanged}, in the emitted environment
$\Gamma$.
\end{definition}
\begin{remark}
Inductive blocks are the one kind of entry the lowering pass copies verbatim: the \lbox{} target
reads only the parameter count and the per-constructor field counts off them, so there is nothing to
lower. This matches MetaCoq's environment erasure, which likewise copies a block essentially
unchanged ([S], Sec. 7.3). \code{IndEmitted} is the state-scoped form of \code{LowerEnv}'s inductive
clause, and it is what Theorem~\ref{thm:RegInvShape-prime-lowerEnv} converts back into the unscoped
clause once the registry is saturated.
\end{remark}
\begin{lemma}[Erasure at the empty context has no free variable]
\label{lem:Erases-noFVar}
\lean{LeanToLambdaBox.Erases.noFVar, LeanToLambdaBox.Erases.noFVar_of_noFVars, LeanToLambdaBox.VLCtx.find?_inr_cons_none}
\leanok
\uses{def:Erases, def:hasFVar, def:lean4lean-grounding}
If $e$ erases to $t$ at a model local context that binds no free variable, then $t$ contains no free
variable node at all; in particular an erasure image taken at the \textbf{empty} local context is
free-variable free.
\end{lemma}
\begin{proof}
\leanok
\uses{def:Erases}
Induction over the erasure derivation. The variable arm is the relation's only producer of a
free-variable node, and its premise that the context answers the lookup contradicts the hypothesis;
the two binder arms extend the context with an \emph{anonymous} de Bruijn entry, which answers no
free-variable lookup --- that is the auxiliary lookup lemma cited above; every other arm is
structural or immediate, the box arm included, since $\Box$ produces no free variable. The
empty-context corollary is the special case where the context answers no lookup at all.
\end{proof}
\begin{remark}
Every producer of \code{RegInvShape'} gets the free-variable-freedom clause this way, the
specification content recording each declared body as an erasure at the empty context.
Free-variable freedom is \emph{independent} of de Bruijn closedness, since a free-variable node is
closed at every level --- the de Bruijn predicate says nothing about locally-nameless variables ---
and it is what the lowering pass's abstraction step consumes. Note the namespace: this is dot
notation on \emph{this} repository's \Erases{} inductive, and the auxiliary lookup lemma is a
repo-local lemma named into an upstream-looking namespace.
\end{remark}
\begin{definition}[The cold-start registry invariant]
\label{def:RegInvShape-prime}
\lean{LeanToLambdaBox.RegInvShape'}
\leanok
\uses{def:IndEmitted, def:SpecContent, def:IndCovered, def:ElimDecl, def:Lower, def:LowerFix, def:Erasure-EraseM, def:Kername, def:LBTerm-envLookup}
$\code{RegInvShape'}\ \code{env}\ \code{bo}\ \Gamma_{\mathrm{spec}}\ s$ relates a \textbf{fixed}
\lbox{} specification environment $\Gamma_{\mathrm{spec}}$ --- the pre-lowering image of the source
environment, the environment MetaCoq's environment erasure talks about --- to the mutable
\code{Erasure.ErasureState} $s$ the shipping eraser has built. Twelve fields in three groups.
\emph{(i) Three clauses about $\Gamma_{\mathrm{spec}}$ alone, fixed for the whole run:}
\code{spec}, its entries say of the source what the specification content demands;
\code{specClosed}, its bodies are de Bruijn closed; \code{specFVarFree}, its bodies mention no free
variable. \emph{(ii) Two coverage clauses:} \code{consts}, every name registered as a constant is
declared in $\Gamma_{\mathrm{spec}}$ at its canonical kername; \code{inds}, every name registered as
an inductive is covered. \emph{(iii) Seven clauses about the emitted environment:} \code{keys}, its
keys are distinct; \code{defs}, a body declared by both environments is either the \Lower{} image of
the specification body or the $j$-th member of a lowered fixpoint block; \code{defsTotal}, every
registered non-runtime-key constant with a specification body is emitted with \emph{some} body;
\code{axioms}, a body-less specification constant is emitted body-less or pruned away;
\code{indsEmitted}, every registered inductive's block is emitted unchanged; \code{sub}, pruning
only removes keys; \code{closed}, all emitted bodies are closed.
\end{definition}
\begin{remark}
Two design decisions are visible in the shape. $\Gamma_{\mathrm{spec}}$ is \textbf{fixed}, not grown
alongside the state, because \Lower{} is \emph{not monotone} in its environment --- its fixpoint arm
carries a negative guard on runtime keys, and declaring a block can turn a key into a runtime key,
invalidating already recorded \Lower{} facts. And the \code{defsTotal} and \code{indsEmitted}
clauses are \textbf{scoped to the registry} rather than to $\Gamma_{\mathrm{spec}}$, so that they
hold vacuously at the empty state and collapse to \code{LowerEnv}'s unscoped clauses only under the
saturation premise of Theorem~\ref{thm:RegInvShape-prime-lowerEnv}. The names \code{specEnv},
\code{lowerEnv}, \code{defns} and \code{erasesEnv} are \emph{not} fields of this structure: they are
theorems in another module (Section~\ref{sec:coldstart:specenv}), and three docstrings in the tree
get this wrong.
\end{remark}
\begin{lemma}[The invariant at the initial state]
\label{lem:RegInvShape-prime-empty}
\lean{LeanToLambdaBox.RegInvShape'.empty}
\leanok
\uses{def:RegInvShape-prime}
At the default \code{ErasureState}, the invariant reduces to its three clauses about
$\Gamma_{\mathrm{spec}}$ alone: given the specification content, closed bodies and free-variable-free
bodies of $\Gamma_{\mathrm{spec}}$, the invariant holds at the empty state.
\end{lemma}
\begin{proof}
\leanok
\uses{def:RegInvShape-prime}
Nine of the twelve clauses are vacuous or trivially true at the empty state: \code{consts},
\code{inds}, \code{defsTotal} and \code{indsEmitted} have unsatisfiable antecedents; \code{keys} is
distinctness of the empty list; \code{defs}, \code{sub} and \code{closed} have no emitted entry to
speak about; \code{axioms} takes its ``pruned'' disjunct. The remaining three are the hypotheses.
\end{proof}
\begin{remark}
This is exactly what makes a \emph{cold} run --- the run \code{Erasure.erase} performs --- provable
at all: the invariant costs nothing at entry.
\end{remark}
\begin{lemma}[Registering a body-less constant]
\label{lem:RegInvShape-prime-addAxiom}
\lean{LeanToLambdaBox.RegInvShape'.addAxiom, LeanToLambdaBox.RegInvShape'.addAxiom_run}
\leanok
\uses{def:RegInvShape-prime, def:Erasure-register_inductive, def:Erasure-addAxiomState}
Registering a body-less constant $n$ preserves the invariant, provided $\Gamma_{\mathrm{spec}}$
declares $n$ body-less too and $n$'s kername does not already occur among the emitted keys. The same
holds of a real monadic run of \code{Erasure.addAxiom}, with no further hypothesis.
\end{lemma}
\begin{proof}
\leanok
\uses{def:RegInvShape-prime, thm:Erasure-run_addAxiom_ok}
The state delta conses a body-less entry at $n$'s kername to the emitted declarations and inserts
$n$ into the constant registry; a one-line unfolding of the delta lets the lookup-across-a-cons
lemmas fire, and a case split on the registry insertion into ``the inserted name or one already
known'' discharges \code{consts} and \code{defsTotal}. The run form is the delta form composed with
the run algebra's \code{addAxiom} lemma, which reports the exact post-state, the \code{panic!}
fall-through included.
\end{proof}
\begin{lemma}[The non-recursive exit preserves the invariant]
\label{lem:RegInvShape-prime-constCons}
\lean{LeanToLambdaBox.RegInvShape'.constCons}
\leanok
\uses{def:RegInvShape-prime, def:Erasure-nonrecConstState, def:Lower, def:LowerFix, def:ElimDecl, def:LBClosed}
Storing a constant $n$ with emitted body $t$ preserves the invariant provided:
$\Gamma_{\mathrm{spec}}$ declares $n$ with some body $b_0$; $t$ is either the \Lower{} image of
$b_0$ or the $j$-th projection of a lowered fixpoint block whose $j$-th key is $n$; $t$ is de Bruijn
closed at level $0$; and $n$'s kername is fresh among the emitted keys.
\end{lemma}
\begin{proof}
\leanok
\uses{def:RegInvShape-prime}
As for the axiom step, with the non-recursive state delta unfolded and the registry insertion split
into the two cases; the \code{defs} clause is the lowering hypothesis at the new key and the
invariant elsewhere, \code{closed} is the closedness hypothesis, and freshness gives \code{keys} and
the two lookup-across-a-cons lemmas.
\end{proof}
\begin{remark}
The workhorse preservation step. It is also the step
Lemma~\ref{lem:RegInvShape-prime-recConst} iterates for block members, and the step the
one-definition cold-run fixture of \code{SpecEnv.lean} uses. Its closedness premise is supplied on a
real run by Theorem~\ref{thm:visitExpr_shape_all}.
\end{remark}
\begin{lemma}[The block exit preserves the invariant]
\label{lem:RegInvShape-prime-recConst}
\lean{LeanToLambdaBox.RegInvShape'.recConst, LeanToLambdaBox.regInvShape'_foldl_recConstStep}
\leanok
\uses{def:RegInvShape-prime, def:Erasure-nonrecConstState, def:LowerBlock}
A whole mutual block is registered in one step and the invariant survives: if the block's
\code{LowerBlock} witness holds at $\Gamma_{\mathrm{spec}}$, the member names sit at the block's own
indices, every member key is fresh against the pre-existing emitted environment, and the member keys
are pairwise distinct, then the invariant holds at the block-registering state. The
member-by-member version iterates the same step along the fold.
\end{lemma}
\begin{proof}
\leanok
\uses{lem:RegInvShape-prime-constCons, def:LowerBlock}
The block-registering state is a fold of a step which is the non-recursive one at a fixpoint-node
body, so each member is an instance of Lemma~\ref{lem:RegInvShape-prime-constCons}: the disjunct for
the \code{defs} clause is supplied by the block's own \code{LowerBlock} witness and closedness by
that witness's closedness projection. Freshness of a member's kername against the entries the
\emph{earlier} members of the same block already consed is exactly the distinctness hypothesis; the
induction is structural on the member list.
\end{proof}
\begin{remark}
The distinctness premise is exactly what the eraser's own obligation
\code{EraserAsks.block\_keys\_distinct} assumes --- a field that is \emph{refuted in general} by a
mutual unsafe definition together with its unsafe-recursive twin (finding F-UNSAFEREC), so every
statement downstream of it is bounded by that.
\end{remark}
\begin{lemma}[A prefix of body-less entries]
\label{lem:RegInvShape-prime-axiomCons}
\lean{LeanToLambdaBox.RegInvShape'.axiomCons, LeanToLambdaBox.regInvShape'_axiomPrefix}
\leanok
\uses{def:RegInvShape-prime, def:LBTerm-envLookup}
Consing one body-less \emph{entry} --- as opposed to registering a \emph{name}: the constant
registry is untouched --- preserves the invariant, provided $\Gamma_{\mathrm{spec}}$ declares that
key body-less and the key is fresh. By induction the same holds of a whole prefix of such entries
with distinct keys.
\end{lemma}
\begin{proof}
\leanok
\uses{def:RegInvShape-prime}
The cons case is the lookup-across-a-cons lemmas at a fresh key, with \code{consts}, \code{inds} and
\code{defsTotal} unchanged because no name was registered. The prefix case is an induction on the
prefix.
\end{proof}
\begin{remark}
This is the shape \code{Erasure.register\_inductive}'s cold branch leaves for each \code{@[extern]}
constructor it axiomatises, without the caller seeing the individual \code{addAxiom} runs. The
prefix form has no consumer outside its own file.
\end{remark}
\begin{lemma}[The emitted block entry]
\label{lem:RegInvShape-prime-blockCons}
\lean{LeanToLambdaBox.RegInvShape'.blockCons}
\leanok
\uses{def:RegInvShape-prime, def:OneInductiveBody, def:Erasure-addAxiomState}
Consing the mutual-inductive block entry that the eraser's registration state emits preserves the
invariant, provided $\Gamma_{\mathrm{spec}}$ holds that same block at that key.
\end{lemma}
\begin{proof}
\leanok
\uses{def:RegInvShape-prime, def:IndEmitted}
The lookup-across-a-cons lemmas at a fresh key; the premise is the unscoped inductive clause read at
one key, so \code{indsEmitted} is preserved and no body clause is touched, a block entry declaring
no body.
\end{proof}
\begin{lemma}[Two transports at a larger state]
\label{lem:RegInvShape-prime-stateCongr}
\lean{LeanToLambdaBox.RegInvShape'.stateCongr, LeanToLambdaBox.RegInvShape'.indsGrow}
\leanok
\uses{def:RegInvShape-prime, def:IndEmitted, def:IndCovered}
The invariant transports across a state whose declarations and inductive registry are unchanged but
whose \emph{constant} registry may have grown, provided every new name is either one the old state
knew or one $\Gamma_{\mathrm{spec}}$ declares body-less. Dually, it transports across a grown
\textbf{inductive} registry, provided every name the registry now knows is covered and has its block
emitted.
\end{lemma}
\begin{proof}
\leanok
\uses{def:RegInvShape-prime}
All clauses not mentioning the moved registry are rewritten across the equalities. In the first
transport the \code{defsTotal} clause is discharged in the new-name case by \emph{vacuity}: a
body-less $\Gamma_{\mathrm{spec}}$ entry does not declare a body, and body-less is the only kind
\code{addAxiom} adds.
\end{proof}
\begin{lemma}[The inductive path, at the run]
\label{lem:RegInvShape-prime-register_inductive_run}
\lean{LeanToLambdaBox.RegInvShape'.register_inductive_run}
\leanok
\uses{def:RegInvShape-prime, def:IndEmitted, def:Erasure-register_inductive, def:Erasure-ConstExt, def:IndCovered}
Let \code{register\_inductive} be run at a name the inductive registry does not yet hold --- the
cold branch --- ending in state $s_1$. Then the invariant holds at $s_1$ \textbf{provided} five
conditions read off the \emph{final} state: the emitted keys are distinct; every body-less final
entry is declared body-less in $\Gamma_{\mathrm{spec}}$; the emitted block at the call's block
kername equals $\Gamma_{\mathrm{spec}}$'s; and each newly registered constant or inductive is either
already present or covered.
\end{lemma}
\begin{proof}
\leanok
\uses{lem:RegInvShape-prime-axiomCons, lem:RegInvShape-prime-stateCongr, lem:RegInvShape-prime-blockCons, thm:Erasure-run_register_inductive_cold_ok}
The intermediate states the cold branch builds --- the \code{@[extern]} axiom prefix, then the block
cons --- are not exposed by the run lemma, so they are rebuilt from the prefix report of the run
algebra's constant-extension record, and the chain axiom prefix, state transport, block cons,
registry growth is applied.
\end{proof}
\begin{remark}
Stated honestly, this lemma \emph{assumes part of its own conclusion}: key distinctness at $s_1$ is
the \code{keys} field of the very invariant it concludes there, plus three further conditions read
off the final state. It is therefore a \emph{transport at a known-good final state}, not a
self-contained preservation step; the module docstring says as much --- key distinctness is not
maintainable unconditionally, because every writer prepends and environment lookup is
first-match-wins.
\end{remark}
\begin{lemma}[A parameter-free body is its own instantiation]
\label{lem:instantiateLevelParams_eq_self}
\lean{LeanToLambdaBox.instantiateLevelParams_eq_self}
\leanok
A source term with no level parameter is unchanged by level instantiation, at any substitution.
\end{lemma}
\begin{proof}
\leanok
\uses{asm:lean-reflection-axioms}
Unfold the instantiation and apply the corresponding core lemma, whose guard is exactly the absence
of a level parameter. The measured footprint of this lemma is the four \code{Lean.Expr}/
\code{Lean.Level} reflection equations and the two \code{bv\_decide} certificates that the
representation lemmas of \code{Lean.Expr} carry, and no \code{sorryAx}.
\end{proof}
\begin{theorem}[Level-scope transport at a parameter-free body]
\label{thm:erases_any_scope_of_paramFree}
\lean{LeanToLambdaBox.erases_any_scope_of_paramFree}
\leanok
\uses{def:Erases, def:NoMaxLevels, def:lean4lean-grounding}
\lbox{} carries no universe levels, so an erasure derivation at one level scope is a derivation at
\emph{every} scope, once the source body has no level parameter left to instantiate and contains no
maximum-level expression: if $b$ has no level parameter, satisfies \code{NoMaxLevels}, and erases to
$b_0$ at the empty local context in scope \code{Us}, then it erases to $b_0$ at the empty local
context in every scope.
\end{theorem}
\begin{proof}
\leanok
\uses{lem:instantiateLevelParams_eq_self, thm:Erases-instL, asm:lean-reflection-axioms}
Instantiate all of \code{Us} at level zero, a well-formed instantiation in any scope, and move the
derivation with the erasure relation's level-instantiation theorem, whose own side condition is
\code{NoMaxLevels}. The instantiation is the identity on a parameter-free body
(Lemma~\ref{lem:instantiateLevelParams_eq_self}) and on the empty local context, so the moved
derivation is the original one read at the new scope. The measured footprint is the same six
reflection names, with no \code{sorryAx}.
\end{proof}
\begin{remark}
This is the one and only bridge from the specification content's delta column, which records
\emph{one} level scope --- the scope a run happened to erase at --- to the environment relation's
delta column, which asks for \emph{every} scope and every instantiation. The gap is real: chapter
\ref{chap:source}'s refutation \code{defns\_needs\_paramFree} refutes the unguarded clause at a body
carrying a level parameter, since the only \Erases{} arm at a sort is the box arm and its
translation witness needs the level in scope. MetaCoq's environment-erasure analogue quantifies over
every universe instantiation because Rocq's constants are universe-polymorphic; the delta column of
the cold-start route is available only at level-parameter-free bodies, and that is a restriction
relative to [S], Sec. 7.3.
\end{remark}

\section{From the invariant to the environment relation}
\label{sec:coldstart:specenv}

The four theorems of this section live in \code{LeanToLambdaBox/SpecEnv.lean}, a module whose other
declarations belong to Chapter~\ref{chap:source}. They are given here because their subject is
\code{RegInvShape'} and their consumer is Theorem~\ref{thm:bridgeEnv_of_regInv}.

\begin{lemma}[The invariant, read as a specification environment]
\label{lem:RegInvShape-prime-specEnv}
\lean{LeanToLambdaBox.RegInvShape'.specEnv}
\leanok
\uses{def:RegInvShape-prime, def:SpecEnv}
The registration invariant is in particular a specification environment for the run state: the four
fields of \code{SpecEnv} are exactly the invariant's specification-side ones.
\end{lemma}
\begin{proof}
\leanok
\uses{def:RegInvShape-prime}
Field by field: the content clause, the two coverage clauses, and free-variable freedom.
\end{proof}
\begin{remark}
The existential wrapper \code{SpecEnv.exists} of Chapter~\ref{chap:source} says that a run which
maintains an invariant \emph{has} a specification environment, and that it is the one the invariant
was maintained against --- derived rather than posited.
\end{remark}
\begin{theorem}[The emitted environment is the lowered image]
\label{thm:RegInvShape-prime-lowerEnv}
\lean{LeanToLambdaBox.RegInvShape'.lowerEnv}
\leanok
\uses{def:RegInvShape-prime, def:LowerEnv, def:RegSaturated}
If the invariant holds at $s$ and the run is \emph{saturated} --- every declared non-runtime-key
definition of $\Gamma_{\mathrm{spec}}$ is some registered constant's and every declared block some
registered inductive's --- then $\LowerEnv{}\ \Gamma_{\mathrm{spec}}\ s.\code{gdecls}$: the emitted
environment is the lowered, pruned image of the specification environment.
\end{theorem}
\begin{proof}
\leanok
\uses{def:RegInvShape-prime, def:RegSaturated}
All eight clauses of \code{LowerEnv} --- \code{keys}, \code{defs}, \code{defsTotal},
\code{axioms}, \code{inds}, \code{sub}, \code{closed}, \code{specClosed} --- come from the
correspondingly named clauses of the invariant; the two \emph{scoped} clauses are unscoped by the
two saturation premises, which produce the registered name behind a declared key.
\end{proof}
\begin{remark}
Without saturation a pruned-away definition would satisfy the scoped clause vacuously, which is
exactly why saturation is a separate premise rather than part of the invariant. The doc comment
records that a former lambda-headedness clause, refuted at every specification environment holding a
non-lambda body, is gone: its content moved to a condition on the blocks the pass \emph{builds}.
\end{remark}
\begin{theorem}[The delta column, from the run's own records]
\label{thm:RegInvShape-prime-defns}
\lean{LeanToLambdaBox.RegInvShape'.defns}
\leanok
\uses{def:RegInvShape-prime, def:Erases, def:ReachableFrom, def:NoMaxLevels}
Assume the invariant, that every kername reachable from a program $t$ is declared in
$\Gamma_{\mathrm{spec}}$, and that every compiler body in the table is level-parameter free with no
maximum-level expression. Then for every tabled constant $c$ whose kername $t$ reaches,
$\Gamma_{\mathrm{spec}}$ declares $c$ with a body $b_0$, and the compiler body erases to $b_0$ at
\textbf{every} level scope and every instantiation.
\end{theorem}
\begin{proof}
\leanok
\uses{def:RegInvShape-prime, lem:instantiateLevelParams_eq_self, thm:erases_any_scope_of_paramFree, def:SpecContent, asm:lean-reflection-axioms}
The reachability premise puts the reached key in the environment; the specification content's delta
clause reads its entry and returns the erasure witness at the one level scope the run erased at; the
level-parameter premise then makes the clause's unguarded quantification reachable --- the
instantiation is the identity (Lemma~\ref{lem:instantiateLevelParams_eq_self}) and the scope moves
freely (Theorem~\ref{thm:erases_any_scope_of_paramFree}).
\end{proof}
\begin{remark}
The names carrying a prime could not be measured directly by the axiom sweep; the footprint is read
off the consumer's row instead. Theorem~\ref{thm:bridgeEnv_of_regInv}, which spends this theorem and
nothing else beyond two projections, is measured at nine names --- the three standard ones plus the
six reflection names of the level-instantiation cluster --- so this theorem carries exactly those
and \textbf{no} \code{sorryAx}. Some notes guess ``expected \code{sorryAx} via \Erases{}''; the
measurement says otherwise, and the blueprint follows the measurement.
\end{remark}
\begin{theorem}[The environment relation of a run]
\label{thm:RegInvShape-prime-erasesEnv}
\lean{LeanToLambdaBox.RegInvShape'.erasesEnv}
\leanok
\uses{def:RegInvShape-prime, def:ErasesEnv, def:ConstOrigin, def:ReachableFrom}
Assume the invariant, the reachability premise, level-parameter freedom of the tabled bodies, and
that every tabled constant is a plain constant of the model. Then
$\ErasesEnv{}\ \code{env}\ \code{bo}\ \Gamma_{\mathrm{spec}}\ t$: the specification environment
erases the source environment along everything the program $t$ reaches.
\end{theorem}
\begin{proof}
\leanok
\uses{lem:RegInvShape-prime-specEnv, thm:RegInvShape-prime-defns, lem:SpecEnv-erasesEnv, asm:lean-reflection-axioms}
Chapter~\ref{chap:source}'s \code{SpecEnv.erasesEnv} takes three premises the state does not record:
the reachability premise, the strengthened delta column, and the table clause. The first is the
hypothesis, the second is Theorem~\ref{thm:RegInvShape-prime-defns}, and the third is the one clause
no registry fact can supply --- at the bridge it is discharged from the reified table.
\end{proof}
\begin{remark}
This is the analogue of MetaCoq's environment erasure ([S], Sec. 7.3) obtained from a \emph{run}'s
records rather than from a global well-formedness assumption.
\end{remark}

\section{The state-side contract of the bridge}
\label{sec:coldstart:bridgeinv}

\begin{definition}[The block's fixvar map]
\label{def:fixvarMap}
\lean{LeanToLambdaBox.fixvarMap, LeanToLambdaBox.fixvarMap_get?_some, LeanToLambdaBox.fixvarMap_ids_subset, LeanToLambdaBox.BridgeInv.fixvars_ids_subset}
\leanok
\uses{def:Erasure-visitMutual, def:Erasure-prepare_erasure}
$\code{fixvarMap}\ \code{nms}\ \code{ids}$ is the hash map \code{Erasure.visitMutual} installs while
erasing a mutual block: each member's source name --- already stripped of its unsafe-recursive
suffix --- mapped to the fresh identifier minted for it. A hit in the map is a hit at one common
index of the two lists; and two equal-length, duplicate-free name lists that assemble the same map
have the same identifiers, so the freshness discipline recorded for the \emph{installed} pair can be
read at any faithful pair.
\end{definition}
\begin{remark}
Four purely combinatorial companions in the same file are left without nodes. The last declaration
of this node is the form the binder steps consume: at any \code{BlockKeyed} pair, each identifier is
reserved by the ambient generator and is not a model-context variable.
\end{remark}
\begin{definition}[A faithful description of the fixvar map]
\label{def:BlockKeyed}
\lean{LeanToLambdaBox.BlockKeyed}
\leanok
\uses{def:fixvarMap, def:Witness-SourceTable, def:Kername, def:Erasure-EraseM}
$\code{BlockKeyed}\ \code{tbl}\ \code{ctx}\ \code{nms}\ \code{ids}$ holds when the reader's fixvar
map \emph{is} $\code{fixvarMap}\ \code{nms}\ \code{ids}$, the two lists have equal length,
\code{nms} is duplicate-free, and the kername map is injective on the \emph{tabled} names --- for
every tabled name, if its kername occurs among the block's kernames then the name itself occurs
among the block's names.
\end{definition}
\begin{remark}
The eraser installs one pair, and the zip admits many others; each conjunct rules out a specific bad
pair. The length condition excludes the appended pair at which the block conjunct is refuted in
Chapter~\ref{chap:refinement}; duplicate-freedom excludes the duplicated pair, whose identifier list
may hold a freshly opened binder; the fourth conjunct is a \emph{separation} restricted to the
tabled names, because unrestricted it is false --- the kername map is not injective, finding
F-KERNAME --- while at the tabled names it is decided by a Boolean checker. The duplicate-freedom
conjunct's shipping counterpart is the refuted \code{EraserAsks.block\_keys\_distinct}.
\end{remark}
\begin{definition}[The refinement conclusion, in either fixvar mode]
\label{def:ErasesLBMode}
\lean{LeanToLambdaBox.ErasesLBMode, LeanToLambdaBox.ErasesLBAltMode, LeanToLambdaBox.ErasesLBMode.ambient, LeanToLambdaBox.ErasesLBMode.block}
\leanok
\uses{def:BlockKeyed, def:ErasesLB, def:ErasesLBFix, def:ErasesLBAlt}
What a sub-run's emitted term is related to, read off the reader's fixvar mode: outside a mutual
block, the emitted term lies in the composite \code{ErasesLB} --- an erasure followed by the
lowering pass; inside one, the block's own constants have been rewritten to its fix variables, which
is the third factor \code{ErasesLBFix} adds, read at the indices of a \code{BlockKeyed} pair. The
alternative-level twin is the same two-mode reading for one \code{case} alternative. The two
projections extract the respective readings.
\end{definition}
\begin{remark}
This is the conclusion shape that lets \emph{one} induction cover both modes; the two projections
are what the two refinement theorems of Chapter~\ref{chap:refinement} return.
\end{remark}
\begin{definition}[The inductive registry is the model's]
\label{def:IndRegistryModelled}
\lean{LeanToLambdaBox.IndRegistryModelled, LeanToLambdaBox.indRegistryModelled_empty}
\leanok
\uses{def:IndInfo, def:IndArity, def:Erasure-EraseM}
Every entry of the run's inductive registry names the block identifier the model declares for that
name, and --- with constructor-argument pruning pinned off --- its argument masks retain every
field. The empty registry is modelled, which is what the cold start needs.
\end{definition}
\begin{remark}
This is the registry twin of \code{Erasure.CanonicalConstants}, and the fact from which the emitted
constructor, projection and case nodes read their identifiers. It is stated at the pinned
configuration's \code{remove\_irrel\_constr\_args := false}.
\end{remark}
\begin{lemma}[Registration keeps the registry modelled]
\label{lem:run_register_inductive_models}
\lean{LeanToLambdaBox.run_register_inductive_models}
\leanok
\uses{def:IndRegistryModelled, def:ConfigPinned, def:ErasureSpec, def:Erasure-register_inductive}
Under the primitive specifications, at a pinned configuration, and given that the block the call is
made at is the one the ambient environment declares at some name, a successful
\code{Erasure.register\_inductive} run leaves the inductive registry modelled.
\end{lemma}
\begin{proof}
\leanok
\uses{def:IndRegistryModelled, asm:ErasureSpec-block_adequate, asm:ErasureSpec-lookup_adequate, asm:ErasureSpec-prim_monotone}
\code{register\_inductive} mints the block kername paired with the member's index for each member of
the block and, with pruning pinned off, an all-keep mask per constructor; the block-adequacy clause
of \code{ErasureSpec} reads both into the model, and the entry-state invariant covers the names the
registry already held.
\end{proof}
\begin{remark}
The statement is guarded by the invariant at the entry state, so a hand-made registry holding a junk
identifier fails the guard rather than the conclusion. It stands under the \code{ErasureSpec}
hypothesis fields.
\end{remark}
\begin{lemma}[Registration only advances the generator]
\label{lem:run_register_inductive_gen}
\lean{LeanToLambdaBox.run_register_inductive_gen}
\leanok
\uses{def:ConfigPinned, def:ErasureSpec, def:Erasure-register_inductive}
Under the primitive specifications and a pinned configuration, a successful
\code{register\_inductive} run only advances the ambient name generator.
\end{lemma}
\begin{proof}
\leanok
\uses{asm:ErasureSpec-prim_monotone, asm:ErasureSpec-lookup_adequate}
The three primitive calls its two loops make are the constant-info query, the environment read and
the log, each bounded by its \code{ErasureSpec} clause; \code{addAxiom} and the registry writes
leave the world token alone; with pruning pinned off the telescope arm is not taken.
\end{proof}
\begin{remark}
This is one of the two bounds Lemma~\ref{lem:runClosedW_gen} takes as a parameter, because it is
declared below that module in the import graph.
\end{remark}
\begin{definition}[The bridge invariant]
\label{def:BridgeInv}
\lean{LeanToLambdaBox.BridgeInv}
\leanok
\uses{def:BlockKeyed, def:IndRegistryModelled, def:Erasure-CanonicalConstants, def:Erasure-EraseM, def:Witness-SourceTable, def:lean4lean-grounding}
$\code{BridgeInv}\ \code{env}\ \code{Us}\ \code{tbl}\ \code{cfg}_0\ \code{gen}\ \code{ctx}\ s\
\Delta$ is what the eraser's reader, its state, the ambient name generator and the modelled local
context agree on at every point of a run. Eight fields: \code{mlc}, the reader's
\code{Lean.LocalContext} is modelled by $\Delta$ through an ambient lean4lean mixed context that is
well formed and matches both sides; \code{lparams}, the reader's level scope \emph{is} the ambient
one; \code{cfg}, the reader's configuration is the one the statement is made at; \code{kfresh} and
\code{reserved}, every model-context variable is reserved both by the kernel's fixed generator and
by the ambient one; \code{fixvars}, the reader is either outside a mutual block or inside one whose
map is described by a \code{BlockKeyed} pair against distinct identifiers, each reserved and none of
them a context variable; \code{canon}, the constant registry is canonical; \code{indcanon}, the
inductive registry is modelled.
\end{definition}
\begin{remark}
This is the state-side contract of the bridge: every one of the eighteen member steps of
Chapter~\ref{chap:refinement} takes and re-establishes it, and Theorem~\ref{thm:erasure_bridge_of_run}
proves it at the entry state by construction. The \code{lparams} field is one of the three measured
obstructions to closing the open item of Section~\ref{sec:coldstart:bridge}: it is stated at a
\emph{fixed} level scope, and \code{Erasure.visitMutual} re-enters a member body under that member's
own \code{levelParams}.
\end{remark}
\begin{lemma}[Three consequences of the modelled context]
\label{lem:BridgeInv-trlctx}
\lean{LeanToLambdaBox.BridgeInv.trlctx, LeanToLambdaBox.BridgeInv.vlctx_wf, LeanToLambdaBox.BridgeInv.noBV}
\leanok
\uses{def:BridgeInv, def:lean4lean-grounding}
From the invariant: the reader's local context corresponds to $\Delta$ under lean4lean's
local-context translation; $\Delta$ is a well-formed model context at the ambient level scope; and
$\Delta$ has no de Bruijn entry.
\end{lemma}
\begin{proof}
\leanok
\uses{def:BridgeInv, asm:lean-reflection-axioms}
All three are read off the mixed-context witness: a mixed context that is well formed and whose two
projections are the reader's context and $\Delta$ yields the correspondence, the well-formedness the
erasure relation's lemmas take, and the absence of de Bruijn entries. The measured footprint adds
one persistent-container reflection axiom.
\end{proof}
\begin{lemma}[The two frame rules]
\label{lem:BridgeInv-mono}
\lean{LeanToLambdaBox.BridgeInv.mono, LeanToLambdaBox.BridgeInv.mono_state}
\leanok
\uses{def:BridgeInv, def:IndRegistryModelled, def:Erasure-RunConcl}
The invariant is monotone in the name generator: it survives advancing the generator. And it
survives state growth: given a canonical run conclusion from $s$ to $s'$ and the fact that the new
state's registry is still modelled, the invariant transports to $s'$.
\end{lemma}
\begin{proof}
\leanok
\uses{def:BridgeInv, def:Erasure-RunConcl}
Reservations survive generator advancement, so the reservation clauses and the fixvar disjunct
transport; canonicity of the constant registry is preserved by the run conclusion, while the
registry model is \emph{not} implied by it --- a run conclusion bounds the registry's growth only
--- and so is taken as a premise, to be supplied by the sub-run's own report.
\end{proof}
\begin{remark}
These are what let the eighteen-motive induction of Chapter~\ref{chap:refinement} compose sub-runs.
\end{remark}
\begin{lemma}[The local-context correspondence across a binder]
\label{lem:TrLCtx-mkLocalDecl}
\lean{LeanToLambdaBox.TrLCtx.mkLocalDecl, LeanToLambdaBox.TrLCtx.mkLetDecl}
\leanok
\uses{def:lean4lean-grounding}
lean4lean's local-context correspondence extends across a lambda binder and across a \code{let}
binder: if the reader's context corresponds to $\Delta$ and the binder's type translates and is a
type in the model, then the context extended with that declaration corresponds to $\Delta$ extended
with the corresponding model entry; with a value and its typing, the \code{let} analogue holds.
\end{lemma}
\begin{proof}
\leanok
\uses{asm:lean-reflection-axioms}
The new entry answers its own lookup and no other, which the file's own \code{Lean.LocalContext}
lookup lemmas supply; the model side is the corresponding model-declaration extension. The
persistent-structure reflection axiom of \code{Lean.PersistentArray} enters the footprint here.
\end{proof}
\begin{remark}
These are repo-local lemmas \emph{named into} an upstream-looking namespace: their real names are
\code{LeanToLambdaBox.TrLCtx.mkLocalDecl} and \code{.mkLetDecl}, and they must be cited qualified.
lean4lean has the ingredients but not the assembled statement.
\end{remark}
\begin{lemma}[The invariant across a binder push]
\label{lem:BridgeInv-mkLocalDecl}
\lean{LeanToLambdaBox.BridgeInv.mkLocalDecl, LeanToLambdaBox.BridgeInv.mkLetDecl}
\leanok
\uses{def:BridgeInv, def:Erasure-mkLambda, def:lean4lean-grounding}
The invariant survives the eraser's two binder-opening helpers: if the binder's type translates and
is a model type, the new variable is not already a model-context variable, it is fresh for the old
generator and reserved by the new one, and the kernel's fixed generator reserves it, then the
invariant holds at the extended reader, the advanced generator and the extended model context.
\end{lemma}
\begin{proof}
\leanok
\uses{def:BridgeInv, lem:TrLCtx-mkLocalDecl, lem:BridgeInv-mono, asm:lean-reflection-axioms}
The mixed-context field is extended by lean4lean's lambda (respectively \code{let}) extension using
Lemma~\ref{lem:TrLCtx-mkLocalDecl}; the freshness bookkeeping is where the two generator hypotheses
do their work --- the new variable is fresh for the old generator, which keeps it apart both from a
block's fix variables and from every entry already in the modelled context --- and the level scope,
the configuration and the two registry clauses are unchanged.
\end{proof}
\begin{remark}
This is the binder step for \code{visitLambda}, \code{visitAlt} and the two eta-expansion families.
The \code{let} version builds the declaration at the default non-dependent binder annotation, which
is the one lean4lean's \code{let} extension records.
\end{remark}
\begin{lemma}[The fragment condition at a subterm]
\label{lem:Supported-subterm}
\lean{LeanToLambdaBox.Supported.subterm, LeanToLambdaBox.Reaches.mono}
\leanok
\uses{def:Supported, def:Witness-SourceTable}
Reachability in the tabled dependency graph is monotone in the constants a term names, so it
transports from a subterm to its container; consequently the fragment condition \code{Supported}
transports to a subterm, given the shape condition the caller inverts out of its own.
\end{lemma}
\begin{proof}
\leanok
\uses{def:Supported}
The bodies clause of \code{Supported} is reachability-scoped, and a subterm names no constant its
container does not, so the clause transports along the monotonicity lemma; the shape clause is the
hypothesis.
\end{proof}
\begin{remark}
This is what the eighteen motives' recursive premises need in order to descend. Both declarations
are repo-local lemmas named into namespaces that look upstream; cite them qualified.
\end{remark}

\section{The erasure bridge, and what remains assumed}
\label{sec:coldstart:bridge}

\begin{theorem}[The erasure half of the capstone, discharged]
\label{thm:erasure_bridge_of_run}
\lean{LeanToLambdaBox.erasure_bridge_of_run}
\leanok
\uses{def:ErasureSpec, def:EraserAsks, def:UpstreamAsks, def:Witness-SourceTableAdequate, def:TableSafe, def:TableBlocks, def:ConfigPinned, def:CompilerBodies, def:Supported, def:lean4lean-grounding, def:Erases, def:Lower, def:SpecEnv, def:Erasure-visitExpr, def:Erasure-prepare_erasure}
Assume the primitive specifications, the eraser's own obligations, the upstream asks, the three
table hypotheses, a pinned configuration, the compiler-body typing, and that the prepared term
\code{pe} lies in the supported fragment and translates into the model. Suppose
\code{prepare\_erasure} took $e$ to \code{pe} from the empty state under the entry reader and
\code{visitExpr} then took \code{pe} to $t$. Then at \textbf{every} specification environment
$\Gamma_{\mathrm{spec}}$ that the run's final state admits there is an intermediate \lbox{} term
$t_0$ with $\Erases{}\ \code{env}\ []\ []\ \code{pe}\ t_0$ and
$\Lower{}\ \Gamma_{\mathrm{spec}}\ t_0\ t$.
\end{theorem}
\begin{proof}
\leanok
\uses{lem:run_prepare_erasure_state, def:BridgeInv, def:IndRegistryModelled, thm:visitExpr_refines_erasesLB, thm:step_visitExpr, lem:step_visitLiteral, lem:step_visitConstructor, lem:step_visitConst, lem:step5, lem:step_visitAppArgs, lem:step_visitProj, lem:step_visitConstApp, lem:step_visitCtorEta, lem:step_visitCases, lem:step_visitAlt, asm:lean4lean-trust, asm:lean-reflection-axioms}
Lemma~\ref{lem:run_prepare_erasure_state} identifies the post-preparation state with the empty one,
so the term walk really starts at the empty state. There the bridge invariant holds \emph{by
construction}: the entry reader carries no fixvar map, so the fixvar clause takes its left disjunct;
the model context is empty, so the freshness clauses and the mixed-context clause are trivial at the
empty mixed context; the level scope and the configuration are the statement's by \code{rfl}; the
constant registry is canonical vacuously; and the inductive registry is modelled by the empty-case
lemma of Definition~\ref{def:IndRegistryModelled}. The refinement theorem of
Chapter~\ref{chap:refinement} is then applied with all eighteen member steps supplied as proved
terms, and its first motive's conclusion read in the ambient mode gives the composite relation,
whose unfolding is the existential in the conclusion.
\end{proof}
\begin{remark}
This is the only half of the composition that is fully proved from the shipping code, and it is the
technical heart of the chapter. Two honest notes. First, the binder \code{hblk}, which assumes
\code{TableBlocks}, is not used in the proof term --- the file sets
\code{set\_option linter.unusedVariables false} and the doc comment says it is consumed at the
install site rather than here; in fact \code{TableBlocks} is consumed by no theorem in the closure,
its sole consumer in the tree having itself no consumer. Second, the theorem's measured footprint
--- taken with the same axiom sweep that produced the ledger, though the theorem is not itself a
ledger subject --- contains \code{sorryAx} together with the reflection axioms, because it routes
the relevance decision through lean4lean's executable checker; it is the immediate dependency of the
capstone, whose committed footprint in \code{test/ledger.expected} is 33 names including
\code{sorryAx}.
\end{remark}
\begin{definition}[The residual, as one named binder]
\label{def:ErasureBridge}
\lean{LeanToLambdaBox.ErasureBridge}
\leanok
\uses{def:ErasesEnv, def:LowerEnv, asm:ErasureBridge-erasesEnv, asm:ErasureBridge-lowerEnv}
$\code{ErasureBridge}\ \code{env}\ \code{bo}\ \Gamma_{\mathrm{spec}}\ \Gamma\ t_0$ bundles exactly
\textbf{two} environment results: \code{erasesEnv}, that $\Gamma_{\mathrm{spec}}$ erases the source
environment along everything $t_0$ reaches; and \code{lowerEnv}, that the emitted environment
$\Gamma$ is the lowered, pruned image of $\Gamma_{\mathrm{spec}}$.
\end{definition}
\begin{remark}
The structure is a definition; its two fields are \emph{assumptions}, owned by
Chapter~\ref{chap:trust} as Assumption~\ref{asm:ErasureBridge-erasesEnv} and
Assumption~\ref{asm:ErasureBridge-lowerEnv}, both unfulfilled at every one of the eight green rungs.
An older design document records a five-field \code{ErasureBridge}; the code has exactly two, the
other three having been retired --- well-formedness became the capstone's own binder, decided per
rung by the Boolean checker of Theorem~\ref{thm:lbWfPeregrine_of_check}; the simulation became a
single application of erasure correctness at the spine; and box-freedom became a lemma on the
first-order constructor-tree shape (Lemma~\ref{lem:noBox_lower_of_foSpine}). \code{ErasesEnv} is this
development's counterpart of MetaCoq's environment erasure ([S], Sec. 7.3), with the extra
Lean-specific layer \code{LowerEnv} beside it for the compilation steps \Lower{} performs.
\end{remark}
\begin{theorem}[The bridge's two environment fields, derived]
\label{thm:bridgeEnv_of_regInv}
\lean{LeanToLambdaBox.bridgeEnv_of_regInv}
\leanok
\uses{def:RegInvShape-prime, def:ErasureBridge, def:RegSaturated, def:SpecEnv, def:LBWfSpec, def:ErasesEnv, def:LowerEnv, def:ReachableFrom, def:ConstOrigin}
Suppose the registration invariant holds at a state \code{sf} against $\Gamma_{\mathrm{spec}}$, the
run is saturated for that environment, every kername the erased term $t_0$ reaches is declared in
$\Gamma_{\mathrm{spec}}$, every tabled body is level-parameter free with no maximum-level
expression, and every tabled constant is a plain constant of the model. Then all four results the
composition needs hold at once: the specification environment for \code{sf}, the well-formedness
$\code{LBWfSpec}\ \Gamma_{\mathrm{spec}}$, and the two \code{ErasureBridge} fields.
\end{theorem}
\begin{proof}
\leanok
\uses{lem:RegInvShape-prime-specEnv, thm:RegInvShape-prime-erasesEnv, thm:RegInvShape-prime-lowerEnv, def:RegInvShape-prime, asm:lean-reflection-axioms}
A four-way conjunction of projections: the specification environment
(Lemma~\ref{lem:RegInvShape-prime-specEnv}); the pair of the invariant's key-distinctness and
specification-closedness clauses, which \emph{is} \code{LBWfSpec}; the environment relation
(Theorem~\ref{thm:RegInvShape-prime-erasesEnv}) at the three side conditions; and the lowered image
(Theorem~\ref{thm:RegInvShape-prime-lowerEnv}) at the saturation premise. The theorem is exactly one
line of Lean. Its measured footprint in \code{test/ledger.expected} is nine names --- the three
standard ones plus four \code{Lean.Expr}/\code{Lean.Level} reflection equations and two
\code{bv\_decide} certificates --- and \textbf{no} \code{sorryAx}.
\end{proof}
\begin{remark}
This is the proved \emph{interface} to the open problem: it fixes precisely what a future
registration-invariant theorem must deliver.
\end{remark}

\subsection*{The open item}

Every theorem of this chapter about the shipping eraser's \emph{runs} is proved, and every theorem
about the registration \emph{invariant} is proved. What is missing is the sentence that joins them:
\textbf{nothing in the repository produces a \code{RegInvShape'} for a state that a run of the
shipping eraser actually reaches.} The only inhabitants of the invariant are
Lemma~\ref{lem:RegInvShape-prime-empty}, at the empty state, and a hand-built one-definition fixture
in \code{SpecEnv.lean}. The consequence is visible in the capstone (Chapter~\ref{chap:capstone}):
\code{shipping\_erase\_correct\_firstorder} carries the hypothesis \code{hbridge}, which asserts for
the run's final state the existence of a specification environment carrying both
\code{ErasureBridge} fields, and it is discharged nowhere, at any of the eight green rungs.
Theorem~\ref{thm:bridgeEnv_of_regInv} is the interface \code{hbridge} would plug into, and
\code{doc/rework/08-REPAIRS-W5.md} states the theorem that would close it, with the content clause
restated so that the erasure witness is \emph{shared} between the specification environment and the
emitted body, read at an environment that \emph{grows with the state} rather than being fixed. The
same document names three measured obstructions: the level scope (\code{BridgeInv} pins a single
scope, while \code{Erasure.visitMutual} re-enters a member under that member's own
\code{levelParams}); alpha-equivalence (a tabled body is pinned to the prepared source body only up
to alpha-equivalence, deliberately, while \Lower{} is not alpha-closed on its source); and the shape
mismatch between the refinement's ``at \emph{every} specification environment'' reading and the
repair's ``one environment, produced'' threading. The preservation kit of
Section~\ref{sec:coldstart:reginv} is exactly what such a theorem would use, and is at this revision
a preparation with no live proof path to the capstone. The dependency graph shows that boundary
rather than hiding it.

\subsection*{Deviations from the paper and caveats}

(i) The applied-versus-block constructor convention is established \emph{on the run}
(Theorem~\ref{thm:visitExpr_ctorSat}, and the applied-form conjunct of
Theorem~\ref{thm:visitExpr_shape_all}), where MetaCoq ships a verified block/applied transformation
([S], Sec. 7.4) --- and the saturation layer has no consumer. (ii) The delta column is available
only at level-parameter-free bodies (Theorem~\ref{thm:erases_any_scope_of_paramFree}), where
MetaCoq's environment erasure quantifies over every universe instantiation; the general clause is
\emph{refuted} in Chapter~\ref{chap:source}. (iii) \Lower{} is a Lean-specific layer with no
MetaCoq counterpart, so \code{ErasureBridge} carries \code{LowerEnv} beside \code{ErasesEnv}.
(iv) Two heartbeat budgets are load-bearing for build reproducibility: four million for the
induction, one million for the block decomposition. (v) Roughly half of \code{ColdStartRun.lean}
--- the whole mutual-block decomposition layer --- and the whole saturation layer are proved but
unconsumed; they are recorded here as machinery, not as steps of the live chain. (vi) No
\code{sorry}, \code{partial}, \code{unsafe}, \code{axiom} or reflection-based decision procedure
occurs in the three cold-start files; the chapter's only trust dependencies are the hypothesis
bundles it takes as parameters and the \code{Lean.Expr} reflection axioms reached through the
level-transport theorem and the local-context lemmas.
